\documentclass[11pt,a4paper]{article}
\usepackage[catalan]{babel}
\usepackage[utf8]{inputenc}
\usepackage[T1]{fontenc}
\usepackage{geometry}
\usepackage{hyperref}
\usepackage{enumitem}
\geometry{margin=2.5cm}

\title{\textbf{README — Documentació de la Pràctica Final}}
\author{Mark Link Cladera \and Antonio Contestí Coll}
\date{}

\begin{document}
\maketitle

\section{Introducció}

Aquest document descriu de manera estructurada el contingut de la entrega de la
\textbf{Pràctica Final}, incloent la memòria escrita, el material audiovisual,
el pòster, els articles científics consultats i els scripts utilitzats per al
desenvolupament, entrenament i avaluació dels models de classificació i detecció
de nòduls pulmonars sobre el conjunt de dades \textit{NODE21}.

L’objectiu d’aquest \texttt{README} és facilitar la comprensió de l’organització
del material entregat i permetre la correcta localització de cada recurs.

\section{Documentació principal}

A la carpeta arrel de la entrega es troben els següents documents principals:

\begin{itemize}[leftmargin=1.5cm]
    \item \textbf{\texttt{AA\_Practica\_Final\_Link\_Contesti.pdf}}  
    Memòria explicativa del treball, on es descriuen els objectius, la metodologia,
    els models utilitzats, els experiments realitzats, els resultats obtinguts i
    les conclusions finals.

    \item \textbf{\texttt{AA\_Practica\_Final\_Link\_Contesti.mp4}}  
    Vídeo explicatiu generat a partir del \textit{notebook} \texttt{lm\_notebook},
    en el qual es presenta de forma resumida l’enfocament del treball, els models
    desenvolupats i els resultats més rellevants.

    \item \textbf{\texttt{Poster\_Practica\_Final\_Link\_Contesti.pdf}}  
    Pòster acadèmic que sintetitza visualment el contingut del treball, incloent
    el plantejament del problema, la metodologia emprada i els principals resultats.
\end{itemize}

\section{Articles científics consultats}

\subsection{\texttt{PAPERS\_CONSULTATS}}

Aquesta carpeta conté tots els articles científics i treballs de recerca
consultats durant el desenvolupament del projecte.  
Tots els documents inclosos en aquesta carpeta apareixen citats com a referències
a la memòria principal (\texttt{AA\_Practica\_Final\_Link\_Contesti.pdf}).

La seva inclusió permet:
\begin{itemize}
    \item Verificar l’estat de l’art utilitzat.
    \item Consultar les fonts originals dels mètodes i arquitectures emprades.
    \item Garantir la traçabilitat acadèmica del treball.
\end{itemize}

\section{Scripts, models i avaluacions}

\subsection{\texttt{SCRIPTS\_MODELS\_I\_AVALUACIONS\_MODELS}}

Aquesta carpeta conté tot el material tècnic utilitzat per al desenvolupament del
treball, organitzat en dos grans blocs: \textbf{classificació} i \textbf{detecció}.

\subsection{Classificació}

La carpeta de classificació inclou els \textit{notebooks} utilitzats per al
preprocessament de dades, l’entrenament dels models i els estudis d’ablació:

\begin{itemize}[leftmargin=1.5cm]
    \item \texttt{conversion\_a\_png.ipynb}
    \item \texttt{node21\_v1\_cs.ipynb}
    \item \texttt{node21\_v1.1\_cs.ipynb}
    \item \texttt{node21\_v1.2\_cs.ipynb}
    \item \texttt{node21\_v1.2ablacio\_cs.ipynb}
\end{itemize}

Així mateix, existeix una subcarpeta \texttt{models} que conté els models
entrenats. Degut a la seva mida, aquests models no s’inclouen directament a la
entrega, sinó que es proporciona un enllaç a Google Drive.

Els models finals són:
\begin{itemize}
    \item \texttt{efficientnet\_b0\_multicanal.pth}
    \item \texttt{efficientnet\_b0\_multicanal\_Unsharp\_bo.pth}
\end{itemize}

\subsection{Detecció}

La carpeta de detecció està estructurada en diversos subdirectoris:

\subsubsection{\texttt{IPYNB}}

Conté els \textit{notebooks} d’entrenament dels diferents models de detecció:

\begin{itemize}
    \item \texttt{node21\_v2.2\_frcnn\_coco.ipynb}
    \item \texttt{node21\_v2.2\_png\_retinanet.ipynb}
    \item \texttt{node21\_v2.3\_frcnn\_vindn\_png.ipynb}
    \item \texttt{node21\_v2.3\_frcnn\_vindn\_attention\_roi.ipynb}
    \item \texttt{node21\_v2.3\_frcnn\_vindn\_png\_attention\_cbam.ipynb}
    \item \texttt{node21\_v2.3\_frcnn\_vindn\_png\_local\_foss.ipynb}
    \item \texttt{yolo\_v1.ipynb}
\end{itemize}

\subsubsection{\texttt{MODELS}}

Inclou els models de detecció entrenats:

\begin{itemize}
    \item \texttt{checkpoint\_coco\_monocanal\_desbalanceado\_epoch\_02\_20251229\_173827.pth}
    \item \texttt{checkpoint\_epoch\_13\_vidn\_local\_foss\_20251231\_152453.pth}
    \item \texttt{checkpoint\_3canal\_frcnn\_vindn\_epoch\_11\_20251231\_003142.pth}
    \item \texttt{checkpoint\_canal\_frcnn\_vindn\_attention\_cbam\_epoch\_16\_20260101\_104415.pth}
    \item \texttt{checkpoint\_retinanet\_epoch\_04\_20260102\_110739.pth}
\end{itemize}

\subsubsection{\texttt{AVALUACIO\_IPYNB}}

\textit{Notebooks} d’avaluació que generen mètriques comparatives en format CSV:

\begin{itemize}
    \item \texttt{evaluate\_node21\_all\_v2.ipynb}
    \item \texttt{evaluate\_node21\_all\_v2yolo.ipynb}
    \item \texttt{evaluate\_node21\_all\_v2retina.ipynb}
    \item \texttt{evaluate\_node21\_all\_vlocal\_foss.ipynb}
\end{itemize}

\subsubsection{\texttt{AVALUACIO\_CSVS}}

Resultats finals de les avaluacions dels models:

\begin{itemize}
    \item \texttt{model\_comparison\_coco.csv}
    \item \texttt{model\_comparison\_retinanet.csv}
    \item \texttt{model\_comparison\_v2.csv}
    \item \texttt{model\_comparison\_v2b.csv}
    \item \texttt{model\_comparison\_vindn.csv}
    \item \texttt{model\_comparison\_vindn\_attention\_cbam.csv}
    \item \texttt{model\_comparison\_vindn\_local\_foss.csv}
\end{itemize}

\section{Observacions finals}

L’estructura de la entrega ha estat dissenyada per garantir la reproductibilitat
del treball, la traçabilitat dels resultats i la claredat en l’organització del
material.  
Els models finals i els resultats presentats a la memòria corresponen directament
als scripts i avaluacions inclosos a les carpetes descrites en aquest document.

\section{Accés en línia al material}

Tot el material corresponent a aquesta Pràctica Final —incloent la memòria,
el vídeo explicatiu, el pòster, els articles científics consultats, els scripts,
els models entrenats i els resultats d’avaluació— es troba disponible de manera
centralitzada i accessible en línia mitjançant Google Drive.

L’enllaç d’accés és el següent:

\begin{center}
\url{https://drive.google.com/drive/folders/1WJeXDtzi0Cx8hRW65Q4TMMt1ho45nHI9}
\end{center}

Aquest repositori permet consultar i descarregar tots els recursos associats
al treball, garantint la traçabilitat, la transparència i la reproductibilitat
dels experiments realitzats.


\end{document}

